\section{Arithmetic geometry of the theta operator}\label{geo:engine}

The arithmetic-geometric argument uses two estimates on the same column pairing. At each finite place a local valuation gives an integral map after a divisor shift. At infinity the chain and its quotient metric give an upper bound on the height of that map. A separate Hermitian-frame argument turns the resulting maps into a lower bound on a target slope. The independent covolume estimate supplies a numerical input to the pairing construction. In the active proof \code{closed\_family\_admissible\_twist\_link\_alt} supplies \code{hA2}; \code{closed\_linkDIVt\_of\_closed\_bridge\_kernel} then supplies \code{hDIV}, which is passed to \code{topDesWin}. These are assembled in \src{src/Development/AutomorphicArithmetic/ThetaCorrespondence/ThetaPairLifts/ThirdPlaceSeries/TopShape.lean}{40}{local estimate and division bridge}.

\noindent For the signed translate mean one may form a labelled direct sum before composing the first edge. Section~\ref{geo:translate} proves the resulting single slope comparison, and Section~\ref{red:translate} substitutes it into the window estimate. \subsection{The operators in the length-\(m\) theta chain}\label{geo:chain}

We use the column-side realization \code{kerRec}, which calls \code{chKerC}, then \code{chKer}. For ordinary edges \code{chKer} takes \code{matK}. For the last two edges it takes the coupled pair \code{cp1K}, \code{cp2K}. The chain is indexed by \(m_2\) ordinary edges followed by those two, with \(m_2+2=m\) by \code{ChIdx.Ok.hm}. \src{src/Formalism/LocalToGlobal/Columns/FiniteSumNonzeroLargeDepth.lean}{993}{exact length} The distribution for the final edge is \code{top}, with the weight at level \(m\). See \src{src/Formalism/LocalToGlobal/Divisibility/DivisionHeightBoundStatement.lean}{943}{active kernel record} and \src{src/Formalism/LocalToGlobal/KernelLattices/ThDistC.lean}{427}{ordinary and coupled kernels}.

Fix an edge \(i\to i+1\), source and target decorated points \(p,p'\) with lattices \(N,N'\), frames \(f,f'\), and indices \(\alpha,\beta\) in their respective level families. Put \(M=\mathrm{homAt}_S(N,N')\), the maps preserving the localized lattices at the specified places, and abbreviate the positive scales of \code{cAof}, \code{cBof} to \(c_A,c_B\). The actual matrix entry is
\begin{equation}\label{geo:chain-entry}
 K_i(p,p')=\sum_{x\in(c_A^{-1}M)/(c_BM)}
              \Xi_i(x)\,\mathcal T_i(N,f,N',f';c_B,x).
\end{equation}
The sum uses classes in the indicated quotient with chosen lifts. The local tests and distribution are descended on the admissible range. Here \(c_A=\prod_P q_P^{A}\), \(c_B=T_u\prod_P q_P^{B}\), and \(S\) consists of the level places and the places used for the row. In the ordinary case \(\Xi_i\) is the level indicator times a product of \(\Xi_{i,P}(f_P'^{-1}xf_P)\). The indicator tests preservation of both the local decorated lattices and their prescribed maps at the level places. Write \(H_P=C_{i+1,P}\cap N_1\), where \(C_{i+1,P}\) is the local \code{Cbe} and \(N_1\) the specified unipotent subgroup. With \(U_P=N_1\cap\operatorname{stab}(\Lambda_{i+1,P})\), the volume ratio used by \code{XiAt} is \(V_{i+1,P}=\mathrm{relIndex}(U_P,H_P)/\mathrm{relIndex}(H_P,U_P)\), with both indices interpreted as in \code{volN}. Let \(d_{i,P}\) denote the number of auxiliary coordinates \code{FC}. \src{src/Formalism/LocalToGlobal/Columns/FiniteSumNonzeroLargeDepth.lean}{542}{relative-index volume} Then the local factor is
\[
 \Xi_{i,P}(z)=V_{i+1,P}^{-1}q_P^{-B d_{i,P}}
   \sum_{\mathbf t\in\mathrm{valsQ}^{d_{i,P}}}
     \Phi_{i,P}(x_1(z,\mathbf t))
       \psi_P(-\mathrm{chArg}(x_1(z,\mathbf t))).
\]
This is a finite weighted Fourier sum. The negative sign and the power of \(q_P^{-B}\) follow by substituting \code{E1at} into \code{XiAt} and expanding \code{E1}. The distribution is zero unless \(c_B\mid c_Ac_B\) and the map belongs to \((c_Ac_B)^{-1}M\). On its support, with \(M_i(N,f)\) the \code{muInvDet} factor for the edge, it is
\[
 \mathcal T_i=
 \frac{M_i(N,f)}{M_i(N',f')}\,
          \mathrm{jetOp}_{i}(J_i).
\]
The coefficients of \(J_i\) are the descents of the complex jet coefficients of the induced theta distribution after multiplication by \(\chi_F^{-1}\). Applied to a polynomial, \(\mathrm{jetOp}_{i}\) first multiplies by the prescribed determinant powers, then substitutes the block variables, contracts their coefficients with \(J_i\), and projects onto the target weight space. This follows by composing the maps in \code{jetMap} and \code{jetOp} in their given order. See \src{src/Formalism/LocalToGlobal/Columns/FiniteSumNonzeroLargeDepth.lean}{636}{lattice quotient sum}, \src{src/Formalism/LocalToGlobal/Columns/FiniteSumNonzeroLargeDepth.lean}{943}{local Fourier normalization}, \src{src/Formalism/LocalToGlobal/Columns/FiniteSumNonzeroLargeDepth.lean}{1110}{scales and local factors}, \src{src/Formalism/LocalToGlobal/KernelLattices/ThDistC.lean}{31}{jet contraction}, and \src{src/Formalism/LocalToGlobal/KernelLattices/ThDistC.lean}{137}{determinant ratio}.

\begin{proposition}[Composition on class representatives]\label{geo:chain-compose}
Let \(F_i\) denote the equivariant forms on level \(i\), with direct summands indexed by its finite label set. Ordinary steps act on these forms by
\[
 (C_i F_i)_\beta(p')=
       \sum_{\alpha}\sum_{c}
          \frac{K_i(p_c,p')\,F_{i,\alpha}(p_c)}{w_c},
\]
where \(p_c\) are representatives of the classes in label \(\alpha\) and \(w_c\) are the source class weights. For the last two steps, the kernel between a point \(p_0\) at level \(m_2\) and a row point \(d\) is
\[
 L(p_0,d)=\sum_c \frac1{w_c}
                \sum_{Y\in\mathcal Y(p_c)}
                   K_{\mathrm{top}}(p_c,Y,d)
                   K_{\mathrm{pre}}(p_0,p_c,Y).
\]
The terminal period on the resulting column \(C(f)=
L\,C_{m_2-1}\cdots C_0\,\mathrm{start}(f)\) and a row \(G\) is
\[
 \mathcal P(C(f),G)
       =\sum_e \frac1{w_e}
            \left\langle \mathrm{br}_u C(f)(d_e),G(d_e)\right\rangle_{K}.
\]
Here \(L\) in the formula for \(C\) denotes the operator of the displayed coupled kernel with its source-class sum, \(d_e\) are row class representatives, and the pairing is the bilinear \code{pairBK}. The map \(\mathrm{br}_u\) is coefficient extraction in the last-column variables followed by projection to the row weight and the scalar
\[
 \prod_\tau t_\tau^{-n_\tau}
        \prod_\tau c_{\mathrm{Br},\tau}\kappa_{\mathrm{Br},\tau},
 \qquad
 t_\tau=\iota_\tau\bigl(\overline{\pi_{T_u}}\bigr).
\]
\end{proposition}
\begin{proof}
For an ordinary step, substitute the class sum \code{ClsKer.appOf} in \code{Mat.opOf}. The normalization is \(w_c^{-1}\) on the source class, and equivariance makes each term independent of the chosen representative. Replacing \(p_c\) by \(g p_c\) contributes \(\rho(g^{-1})\) on the right of the kernel and \(\rho(g)\) on the form. For the coupled pair the two kernel entries use the \code{XiCpl1}, \code{XiCpl2} tests respectively, on the same admissible coset \(Y\). They are zero off the admissible set. Upon changing the intermediate representative, \(Y\) is transported bijectively and the intervening factors \(\rho(g^{-1})\rho(g)\) cancel in the product. Thus there is one intermediate class weight and one sum over \(Y\), rather than two unrelated coset sums. The recursion for \code{colTo} starts at the identity and composes each ordinary step on the left. The \code{colCpl} formula then applies the coupled pair. Finally \code{PtilFF} evaluates the branching map and pairs at the same row representative, with the row class weight. Substituting \code{brAt} and \code{brStdP} gives the displayed scalar. The formula is a dual pairing against the row input. It requires no identification of either of the last two maps with the Hermitian adjoint of the other.
\end{proof}
\noindent
\src{src/Formalism/LocalToGlobal/Columns/RhoEqOfSmulEq.lean}{372}{class action and normalization},
\src{src/Formalism/LocalToGlobal/Columns/RhoEqOfSmulEq.lean}{679}{matrix action},
\src{src/Formalism/LocalToGlobal/Columns/FiniteSumNonzeroLargeDepth.lean}{263}{chain recursion},
\src{src/Formalism/LocalToGlobal/Columns/FiniteSumNonzeroLargeDepth.lean}{298}{coupled contraction},
\src{src/Formalism/LocalToGlobal/Columns/FiniteSumNonzeroLargeDepth.lean}{128}{branching coefficient},
\src{src/Formalism/LocalToGlobal/Columns/FiniteSumNonzeroLargeDepth.lean}{215}{terminal row pairing}.

The active coordinate comparison \code{pairing\_eq\_colPairing} identifies the rational map \(B_k(u)\) with the column pairing. It uses the period identity supplied by the realized column data and its vanishing on the right coordinate kernel. It then evaluates on a right inverse of the surjective quotient-coordinate matrix to prove the equality on every target coordinate. The finite-place comparison \code{perToVal2F} carries the local period estimate to this rational pairing. Evaluating individual lifts would give neither this descended functional nor its accumulated finite-place constraint. In the ordinary local model the valuation is taken on the actual chain end, using the recursion \(C_{[0]}=1,\ C_{[j+1]}=C_j C_{[j]}\) and the source's compensated generator estimate. At the endpoint its scalar budget contains \(e_w\sum_{j=1}^{m}(g_S(p_w,j)+F(j,w)-E(u,j,w))\), while the period cost subtracts the stabilizer contribution once, as in \eqref{geo:budget} and Lemma~\ref{geo:generators}. These are quantitative bounds on the composition, with the model and comparison hypotheses stated there. At \(p_w=T_u\) the per-layer budget instead aggregates to \(-m\,\mathrm{EatT}(u,w)\) before multiplying by \(e_w\). See \src{src/Development/AbelianVarietyTheory/Proofs/IwasawaDecomposition/RationalPairingEqColumnPairing.lean}{60}{active endpoint comparison}, \src{src/Development/AbelianVarietyTheory/Proofs/IwasawaDecomposition/RationalPairingEqColumnPairing.lean}{23}{rational pairing on the quotient}, \src{src/Formalism/LocalToGlobal/GroupLattices/ServJointF.lean}{42}{finite-place comparison}, \src{src/Formalism/SharpRestriction/GrowthBounds/DataGChain.lean}{85}{local chain recursion}, and \src{src/Formalism/SharpRestriction/Divisibility/WeightDegreeSizes.lean}{887}{level-prime length factor}.


\subsection{Finite-place integrality and divisor translation}

Fix a word \(u\) and a finite place \(w\) of the coefficient field \(K\), with residue characteristic \(p_w\) and ramification index \(e_w=v_w(p_w)\). Use the convention \(v_w(0)=+\infty\) and normalize the other values in \(\mathbb Z\). Let \(m\) be the chain length, \(T_u\) the level prime, and \(\Gamma_u\) the positive integer \code{DGam} divisible by every class stabilizer order in each layer \(0\le i\le m\). \src{src/Formalism/LocalRingCardinality/DataKernels.lean}{25}{common multiple of stabilizer orders} Put
\begin{equation}\label{geo:budget}
 d_w=v_{p_w}(\Gamma_u),\qquad
 s_{u,w}=\sum_{j=1}^{m}\bigl(g_D(j,w)+F(j,w)-E(u,j,w)\bigr).
\end{equation}
Here \(g_D(j,w)\) is the recorded gain with the gain at every dropped prime set to zero. The model gain \(g_S(p_w,j)\) uses the positive parts of the orbital gains. In particular an empty set of selected orbits can leave a nonzero contribution from the complementary orbits. The comparison \(g_D\le g_S\) uses the filled parameters and the inclusion of the exceptional \code{IsCs} primes among the dropped primes. At a dropped prime or a place outside the gain support the recorded gain vanishes. Failed inertia tests and the two generator primes give zero by the local gain equations. In the remaining branch the filled parameters identify the recorded gain with the model before the exceptional-prime mask, and the exceptional branch is ruled out by the drop condition. The nonnegativity of \(g_S\) then completes the comparison. See \src{src/Formalism/SharpRestriction/GrowthBounds/WeilJacquetShalikaLaw.lean}{538}{comparison of recorded and model gains}.

\begin{lemma}[Generator and period budgets]\label{geo:generators}
Let \(A\) be a discrete valuation ring with fraction field \(L\), normalized valuation \(v\), and let \(M\) be an \(A\)-submodule of an \(L\)-vector space. Let \(C\) be an \(L\)-linear chain map. Suppose an input lattice is generated over \(A\) by \(a_i x_i\), with \(C x_i\in \pi^{b_i} M\), \(b_i\in\mathbb Z\), and \(v(a_i)+b_i\ge b\) for every \(i\), where \(b\in\mathbb Z\). Then every input \(f\) satisfies \(C f\in\pi^b M\). If a terminal \(L\)-linear period \(\mathcal P\) satisfies \(\mathcal P(M)\subset\pi^{-d}A\) for an integer \(d\), then \(v(\mathcal P(Cf))\ge b-d\).
\end{lemma}
\begin{proof}
For each generator, \(a_i\pi^{b_i-b}\in A\). Hence \(\pi^{-b}C(a_i x_i)\in M\). The coefficients of any finite linear combination of generators belong to \(A\). Applying \(\pi^{-b}C\) to such a combination therefore gives an element of \(M\), including for the zero input or a cancelling combination. Write \(C f=\pi^b z\) with \(z\in M\). Then \(\pi^{d-b}\mathcal P(C f)=\pi^d\mathcal P(z)\in A\). Negative values of \(b\) and \(d\) are allowed, since the powers are taken in \(L\).
\end{proof}

In the active ordinary-place branch the twisted lattice is generated by products \(a_{\mathcal V'}F_0\). Here \(\mathcal V'\) ranges over subsets of the selected orbits \(\mathcal V\). The complement contributes the \(p_w\)-power budget
\[
 \sum_{t\in\mathcal V\setminus\mathcal V'}
       \max\bigl(0,\lceil \mathrm{Sv}(t)\rceil\bigr).
\]
The contract \code{unitBoundAtTwistedPlaces} bounds the chain on \(F_0\). The orbit inequality \code{bg\_gain\_le} converts the generator-specific bound after compensation into the uniform bound. The Lean proof writes membership using the uniformizer and applies span induction, precisely the mechanism of Lemma~\ref{geo:generators}. It starts from a chain-end estimate on specified generators. The terminal column comparison includes the ordinary and twisted coordinates. When the selected-model branch is unavailable, the comparison instead uses the supplied unselected-orbit lattice bound. See \src{src/Formalism/SharpRestriction/GrowthBounds/AdmissibleIntegers/StatementsOfRemainingBound.lean}{111}{generator compensation and span induction} and \src{src/Formalism/SharpRestriction/GrowthBounds/AdmissibleIntegers/StatementsOfRemainingBound.lean}{162}{ordinary-place model alternatives}.

The resulting column has scalar valuation at least
\[
 e_w\left(d_w+\sum_{j=1}^m(g_S(p_w,j)+F(j,w)-E(u,j,w))\right).
\]
The terminal period on the remaining vector costs at most \(e_w d_w\). Indeed its two lattice-coordinate versions use the integrality at the ordinary or twisted points, respectively, and the loss from the same common multiple. The valuation of the product is therefore at least \(e_w s_{u,w}\), using \(g_D\le g_S\). The loss \(e_w d_w\) occurs once. This is the multiplication in \src{src/Formalism/SharpRestriction/GrowthBounds/WeilJacquetShalikaLaw.lean}{401}{ordinary column and terminal period}. At \(p_w=T_u\) the level-lattice calculation yields the budget \(-m\,\mathrm{EatT}(u,w)\) before multiplication by \(e_w\), and \(d_w=0\) by coprimality. At a generator prime distinct from \(T_u\), one multiplies the test vector by a local unit to bring it into the lattice on which the local bound holds. Dividing by that unit preserves its valuation. These branches and the final integer bound \(\lfloor e_ws_{u,w}\rfloor\) occur in \src{src/Formalism/LocalToGlobal/GroupLattices/AdmissibleAtGeneratorPlaces.lean}{214}{three finite-place cases}. The family version takes a common threshold for the inputs before supplying \code{adm2}. See \src{src/Development/AutomorphicArithmetic/ThetaCorrespondence/ThetaPairLifts/ThirdPlaceSeries/ClosedFamilyAdmissibleTwistLink.lean}{24}{family estimate}. The local generator estimates and the ordinary and twisted comparison maps are the geometric premises in this step.

\begin{proposition}[Integral shifted dual map]\label{geo:shift}
Let \(B:N_K\times R_K\to K\) be bilinear on the scalar extensions of full \(\mathcal O_K\)-lattices \(N,R\). At every finite \(w\), suppose the available period coordinates \(A_w\)-span \(R_w\), where \(A_w=\mathcal O_{K,w}\), and the period equals the evaluation of \(B\) at those coordinates up to a unit. Assume the valuation of the period for every element of \(N_w\) and every such coordinate is at least \(\sigma_w\in\mathbb R\). If \(\tau_w\le\lceil\sigma_w\rceil\) are integers of finite support, the shifted source lattice with local model \(\pi_w^{-\tau_w}N_w\) maps into \(R^\vee\), the integral linear dual of \(R\).
\end{proposition}
\begin{proof}
An integral valuation at least \(\sigma_w\) is at least \(\lceil\sigma_w\rceil\). The ideal \(\pi_w^{\lceil\sigma_w\rceil}A_w\) is an \(A_w\)-module, so comparison by a unit and extension from the spanning coordinates preserve this conclusion. For \(f\in\pi_w^{-\tau_w}N_w\) and \(r\in R_w\), bilinearity now gives \(v_w(B(f,r))\ge\lceil\sigma_w\rceil-\tau_w\ge0\). For a global element of the shifted source and a global element of \(R\), the value thus belongs to every \(A_w\). The intersection of the \(A_w\) inside \(K\) is \(\mathcal O_K\), since a principal fractional ideal with no negative prime exponent is integral. This proves the dual-lattice inclusion.
\end{proof}

The coordinate comparison and spanning premises are supplied as part of the division package. The actual construction introduces them in \src{src/Formalism/LocalToGlobal/Divisibility/RecBridge.lean}{633}{feeding the period estimate into the division bridge}. See also \src{src/Formalism/LocalToGlobal/GroupLattices/ServJointF.lean}{84}{local column comparison} and \src{src/Formalism/QuotientMetricDivisors/HermitianTwistDivisors/TwistedLatticeNaturality.lean}{720}{integrality after the shift}.

\begin{corollary}[Determinant cost of the shift]\label{geo:shift-index}
Suppose the lattices in Proposition~\ref{geo:shift} have a fixed nonsingular infinite-place metric, with \(K\) totally imaginary and \(N\) of rank \(d\). Write \(N'\) for the shifted source and \(l_w=\log|\mathcal O_K/\mathfrak w|/[K:\mathbb Q]\). Then
\[
 \mathrm{lcv}(N')-\mathrm{lcv}(N)=-2d\sum_w l_w\tau_w,\qquad
 \mu_{\rm def}(N')-\mu_{\rm def}(N)=2\sum_w l_w\tau_w\quad(d>0).
\]
For a rational prime \(p\), \(\sum_w l_w v_w(p)=\log p\).
\end{corollary}
\begin{proof}
At a finite place, take an \(A_w\)-basis of \(N_w\). The shifted basis is obtained by multiplying each of its \(d\) vectors by \(\pi_w^{-\tau_w}\), so the determinant has valuation \(-d\tau_w\), including the zero-rank determinant \(1\). The fixed-metric comparison of log-covolumes is twice the sum of determinant valuations with weights \(l_w\). Its factor two is part of the normalization in \src{src/Formalism/QuotientMetricDivisors/HermitianTwistDivisors/TwistedLatticeNaturality.lean}{145}{determinant comparison and slope normalization}. For \(d>0\), \(\mu_{\rm def}=-\mathrm{lcv}/d\), giving the stated sign. Finally \(v_w(p)=0\) away from the places over \(p\), where \(v_w(p)=e_w\) and \(\log|\mathcal O_K/\mathfrak w|=f_w\log p\). The degree formula \(\sum_{w\mid p}e_w f_w=[K:\mathbb Q]\) yields the rational-prime identity. Equivalently, specialize the active valuative product formula to \(p\). All complex embeddings have modulus \(p\), and their squared normalization uses \(2\,\#S_\infty=[K:\mathbb Q]\). \src{src/Formalism/QuotientMetricDivisors/HermitianTwistDivisors/ProdFormula.lean}{39}{valuative product formula}
\end{proof}


\subsection{Archimedean coercivity and slope transfer}

\begin{lemma}[Height of a descended pairing]\label{geo:descent}
Suppose a Hermitian quotient used by the right input of \(B_u\) has a linear section taking that input to a lift of norm at most \(\beta^{-1/2}\) times its norm, with \(\beta>0\). Suppose the chain maps \(C_j\) have operator norms squared at most \(e^{k_j}\), and the terminal pairing has operator norm squared at most \(e^{h_P}\). Then
\[
 \|B_u\|^2\le e^{h_P+\sum_j k_j}/\beta.
\]
In particular, if \(\beta=\kappa\min_a\prod_j\lambda_{j,a}\) with \(\kappa>0\) and all the \(\lambda_{j,a}>0\), the logarithm for a nonzero \(B_u\) is bounded by \(h_P+\sum_j k_j-\log\kappa-\min_a\sum_j\log\lambda_{j,a}\).
\end{lemma}
\begin{proof}
Apply the quotient section to the right input and the chain to the left input. Successive operator-norm bounds give \(\|C_m\cdots C_1 x\|^2\le e^{\sum_j k_j}\|x\|^2\). Applying the terminal pairing and then the section bound proves the claimed norm inequality for all input pairs. If a Hermitian Gram operator \(G\) controls the section norm through \(\|y\|^2=\langle G\widetilde y,\widetilde y\rangle\) on the lifting complement and all its eigenvalues are at least \(\beta\), an orthonormal eigenbasis gives \(\beta\|\widetilde y\|^2\le\|y\|^2\), including \(\widetilde y=0\). In the application the product lower bound for \(\beta\) comes from the common block description of the quotient Gram operator. The zero map satisfies the norm inequality and needs no logarithm.
\end{proof}

Averaging the logarithmic estimate over infinite places yields \(\mathrm{htMap}(B_u)\le h(u)\) for the nonzero maps with the chosen word height budget. The normalization uses the mean of \(\log\|B_u\|_v^2\). See \src{src/Formalism/QuotientMetricDivisors/HermitianTwistDivisors/DescendedNorm.lean}{823}{chain in the quotient metric} and \src{src/Formalism/LocalToGlobal/Divisibility/ConstantGeTwentyEight.lean}{123}{average height}. The companion estimate on target slope plus word height also needs the border-minor bound on the base side and the row-to-column comparison in \code{P.extra.passage}. These inputs are separate from the bound on the operator norm. They are combined in \src{src/Formalism/LocalToGlobal/Divisibility/ConstantGeTwentyEight.lean}{210}{target-slope budget}. The actual bridge packages the integral inclusion, the height bound for nonzero maps, and this combined estimate as \code{WordFacts}. See \src{src/Formalism/LocalToGlobal/Divisibility/RecBridge.lean}{128}{wordwise division data}.

\begin{lemma}[Coercivity of a nearly tight frame]\label{geo:coercivity}
Let \(A_i:V_i\to \mathcal H\) be a finite family of maps of finite-dimensional complex Hilbert spaces. Put \(r=\sum_i\dim V_i>0\), \(D=\dim\mathcal H\), and \(Q=\sum_i A_iA_i^*\). Suppose
\[
 \sum_i\|A_i^*A_i-I_{V_i}\|_{\mathrm{HS}}^2\le er,\quad
 0\le e\le\tfrac14,\qquad
 \operatorname{tr}(Q^2)\le (1+\eta)(\operatorname{tr}Q)^2/D,\quad \eta\ge0.
\]
The second inequality is imposed on the nonzero ambient space. For subspaces \(U_i\subseteq V_i\), write \(q=\sum_i\operatorname{codim}_{V_i}U_i/r\). Then
\[
 \frac{\operatorname{codim}_{\mathcal H}\sum_i A_iU_i}{D}
       \le\eta+4q+4\sqrt e.
\]
Positive frame weights are allowed by putting their square roots into \(A_i\).
\end{lemma}
\begin{proof}
Write \(E_i=A_i^*A_i-I\). Then
\[
 \operatorname{tr}Q=r+\sum_i\operatorname{tr}E_i
 \ge r-\sqrt{r\sum_i\|E_i\|_{\mathrm{HS}}^2}
 \ge r(1-\sqrt e).
\]
This also rules out a zero ambient space using just the first hypothesis. Let \(P\) project onto \(Z=(\sum_i A_iU_i)^\perp\) and set \(d=\dim Z\). Since \(P A_i\) kills \(U_i\), if \(S_i\) projects onto \(U_i^\perp\) then
\[
 \operatorname{tr}(PQ)\le\sum_i\operatorname{tr}(S_iA_i^*A_i)
 \le rq+\sqrt{rq\sum_i\|E_i\|_{\mathrm{HS}}^2}
 \le r(q+\sqrt{qe}).
\]
For the first inequality, take an orthonormal basis adapted to \(U_i\) and bound the squared norms of the projected images by the squared norms of the images. With \(M=\operatorname{tr}Q>0\) the second frame hypothesis gives \(\|Q-(M/D)I\|_{\mathrm{HS}}^2\le\eta M^2/D\). Pairing with \(P\) yields \(\operatorname{tr}(PQ)/M\ge z-\sqrt{\eta z}\) for \(z=d/D\). Since \(2\sqrt{\eta z}\le\eta+z\), this lower bound is at least \((z-\eta)/2\). The upper bound is at most \(2(q+\sqrt e)\), using \(q\le1\) and \(1-\sqrt e\ge 1/2\). The conclusion follows. The proof allows \(\dim V_i=0\), \(Z=0\), \(q=0\), and \(e=0\).
\end{proof}

The active linear-algebra estimate uses the trace of the complementary projection in \src{src/Formalism/HarderNarasimhan/HarderNarasimhanProofs/FrameLinearAlgebraFirst.lean}{149}{complementary trace bounds}. In its compressed version \(\mathcal H=\mathcal H_0\oplus\mathcal H_1\), where the target is \(\mathcal H_0\) of dimension \(D_0\ge(1-\delta)D\). If \(M_0\) is the projection to \(\mathcal H_0\) of \(M=\sum_i A_iU_i\), then \(\mathcal H/M\) surjects onto \(\mathcal H_0/M_0\). Division by \(D_0\) gives the preceding relative deficit divided by \(1-\delta\), with \(D_0>0\) since \(0\le\delta<1\). The main application supplies the two frame parameters by compressing the full matrices along isometric embeddings. It substitutes \(\eta=5\sigma+3\sqrt{d_{\rm rk}}\) and \(e=2e_0\) after the rank and Hilbert--Schmidt estimates. See \src{src/Formalism/SizeWindowBrackets/ChainSlopeWindows/JointDModelAppliesUnitLattice.lean}{164}{frame compression estimates} and \src{src/Formalism/SharpRestriction/GlobalSchurBounds/TopDensityNGeSixteen.lean}{538}{actual frame and rank data}.

\begin{theorem}[Finite-layer HN estimate]\label{geo:hn}
Use normalized Hermitian slopes \(\lambda=2\lambda_{\rm Euc}-\log 2+\log|\operatorname{disc}K|/[K:\mathbb Q]\) for a totally imaginary \(K\). Suppose a finite family of lattice maps \(B_i:N_i\to R_*\) of height at most \(h_i\) for nonzero maps has positive total source rank. Fix \(0\le\delta<1\). Assume the target lattice has positive rank \(D_0\) over \(K\), and that at a selected infinite place its Hermitian space identifies with a subspace \(\mathcal H_0\subseteq\mathcal H\) of dimension \(D_0\ge(1-\delta)D\) in the frame of Lemma~\ref{geo:coercivity}. Assume the matrices of the lattice maps at that place are the orthogonal projections of the frame maps to \(\mathcal H_0\), up to multiplication of each matrix by a nonzero scalar. Let \(s\) be the rank-weighted mean of all source slopes after subtracting \(h_i\), and assume these slopes are at most \(s+u\) with \(u\ge0\). Assume every target slope is at least \(s-H\), \(H\ge0\). Then
\begin{equation}\label{geo:hn-result}
 s-\mu(R_*)\le
 \frac{4u\log(1+H/u)+(\eta+4\sqrt e)H}{1-\delta}.
\end{equation}
The first term in the numerator is defined as zero when \(u=0\).
\end{theorem}
\begin{proof}
For \(t\in[0,H]\), truncate each source at adjusted slope \(s-t\). The slope inequality for a morphism of Hermitian lattices gives, for a nonzero image of the corresponding source truncation, minimum slope at least \(s-t+h_i-\mathrm{htMap}(B_i)\ge s-t\). It uses the image lattice with the induced target metric. The image is therefore contained in the target high-slope truncation by HN functoriality at the chosen threshold. For the zero image containment holds directly. The normalized slopes and height account for the factor of two in the Euclidean scale. The slope inequality used here is \src{src/Formalism/HarderNarasimhan/HarderNarasimhanProofs/SlopeIneqMean.lean}{249}{image slope under a height bound}. The truncation containment, including the zero branch, occurs in \src{src/Formalism/HarderNarasimhan/HarderNarasimhanProofs/TheoremFSharp.lean}{277}{images of high-slope lattices}.

Let \(q_t\) be the source fraction below \(s-t\). Since the adjusted source mean equals \(s\),
\[
 s\le q_t(s-t)+(1-q_t)(s+u).
\]
For \(u+t>0\) this gives \(q_t\le u/(u+t)\). When \(u=0\), the mean condition and the common upper bound imply every source slope equals \(s\), so \(q_t=0\) even at \(t=0\). The target fraction strictly below \(s-t\) is no greater than the codimension of the sum of the projected truncated images divided by \(D_0\). The preceding frame bounds thus majorize this fraction by \((\eta+4q_t+4\sqrt e)/(1-\delta)\).

Each target slope \(\lambda\) satisfies \((s-\lambda)_+=\int_0^H\mathbf1_{\lambda<s-t}\,dt\). Equality at a threshold changes no integral. Summing this finite identity gives the layer-cake bound on \(s-\mu(R_*)\). For \(u>0\), the integral of \(u/(u+t)\) is \(u\log(1+H/u)\). For \(u=0\) its replacement \(q_t\) is identically zero. At \(H=0\) the target minimum hypothesis already gives \(s-\mu(R_*)\le0\). These are the rank-count and integration steps in \src{src/Formalism/HarderNarasimhan/HarderNarasimhanProofs/TheoremFSharp.lean}{308}{frame rank deficit} and \src{src/Formalism/HarderNarasimhan/HarderNarasimhanProofs/TheoremFSharp.lean}{338}{finite-slope layer cake}.
\end{proof}

The application takes the weakened bound with an added \(\log|I|\) for the nonempty index set of words. The division estimate bounds \(s\) below by the target slope plus the gain, minus the numerical losses. Applying Theorem~\ref{geo:hn} bounds the same \(s\) above by the target slope plus the frame and counting losses. Subtracting eliminates the target slope, giving the one-sided gain estimate entering \(\mathrm{ASK}\). The height and slope transfer is \src{src/Formalism/SharpRestriction/GlobalSchurBounds/QuadraticLogDerivCongr.lean}{699}{conversion to the actual frame data}. See \src{src/Formalism/SizeWindowBrackets/ChainSlopeWindows/LevelOfCongruence.lean}{50}{lower bound on the source mean} and \src{src/Formalism/SizeWindowBrackets/ChainSlopeWindows/LevelOfCongruence.lean}{113}{slope elimination}.

\subsection{Matching theta integrals and constructing the evaluation lattice}

In the tensor Weil model for an admissible Hermitian triple \(V,V',W\), write
\[
 \Theta_\omega(g,h;\Phi)=\operatorname{thetaSum}_F(\omega(g,h)\Phi),
 \qquad I_V(\Phi)(h)=
 \int_{\mathcal F_V}\Theta_\omega(g,h;\Phi)\,d\mu_V(g).
\]
Here \(h\in U(W)(\mathbb A_F)\) and \(\mathcal F_V\) is a fundamental domain for \(U(V)(F)\backslash U(V)(\mathbb A_F)\). The corresponding domain and measure on the other side are those of its Weil model. In particular the global comparison constant depends on the fixed matching maps and model measures. The volume-one normalization below applies to the unipotent constant term. The definitions are in \src{src/Development/AbelianVarietyTheory/BlockTheory/DivisorBase/WeilTensorModels.lean}{286}{theta series and integrals} and \src{src/Development/AbelianVarietyTheory/Definitions/Divisors/CTSmallCone.lean}{82}{normalized constant term}.

At good finite places the two spherical vectors have the same nonzero image in the common quotient. For a finite set \(S\) containing the bad places, polynomial matching of pure tensors identifies the tensor of the real image with the images in the finite-place quotients:
\[
 R_\infty(\phi_\infty)\otimes\bigotimes_{v\in S}R_v(\phi_v)
 =
 R'_\infty(\phi'_\infty)\otimes\bigotimes_{v\in S}R'_v(\phi'_v).
\]
The condition extends linearly to finite sums. At good places the relevant kernel is the span of the differences \(\omega(g,1)\phi-\phi\), so a linear functional invariant under \(g\) factors through the quotient because the functional vanishes on each difference and hence on the kernel. The general matching condition is a limit of polynomial matching, with a support condition on the right in the split matrix model that forces vanishing on rank-deficient matrices. See \src{src/Development/AbelianVarietyTheory/BlockTheory/DivisorBase/WeilTensorModels.lean}{459}{polynomial matching and matched limits}.

\begin{proposition}[Proportionality from the matching inputs]\label{geo:theta}
In the tensor Weil models, fix real and finite matching maps satisfying \code{IsMatched}. Assume the anisotropic bottom comparison, the constant-term descent and reconstruction identities, nonvanishing of the descent factors, a nonzero test at the bottom, and continuity along the matched approximations. Assume also that the right-hand polynomial test obtained by nonvanishing admits a left-hand match. Then there is a unique \(c\in\mathbb C^\times\) such that \(I_V(\Phi)(h)=cI_{V'}(\Phi')(h)\) for every matched pair and every \(h\), and some matched pair has nonzero right-hand integral.
\end{proposition}
\begin{proof}
First use polynomial vectors and induct on the sum \(p+q\) of the two Witt indices in the comparison data. The bottom comparison supplies the step \(p=q=0\). When \(q>0\), remove an adapted hyperbolic plane in the primed tower and the corresponding plane on the target. If \(N\) is the resulting unipotent group, take the constant term \(\int_{\mathcal F_N} f(nh)\,d\mu_N(n)\), where \(\mathcal F_N\) represents \(N(F)\backslash N(\mathbb A_F)\) and has measure one. Its two summands are identified by the specified Levi transformation rules. For \(a'=b+2\), set \(k=\lfloor(a-(b+2))/2\rfloor\) using natural subtraction, write \(c_1\) for the comparison constant with Witt index sum lower by one, and set \(F_0=\mathrm{stepFac}(-k)\), \(L_0=L^{S_0}(\nu,k+1)\) in the supplied descent identities. With the nonzero local constant \(C\), the compatibility identity to feed the reconstruction is \(cL_0=c_1/(CF_0)\). Indeed \(c=c_1/(CF_0L_0)\) exists and is nonzero, and multiplying by the nonzero product verifies the identity. In the branch \(q=0<p\), the unprimed descent gives the comparison in the other direction. The specified uniqueness of the constant-term splitting identifies the first component with the lower-tower component. The nonzero descent factors transport the lower comparison to the upper reconstruction. This is the induction implemented in \src{src/Development/WeilBound/ThetaLifts/ConeHypotheses/CtSmallCone.lean}{88}{two constant-term descents}. The explicit primed coefficient is constructed at \src{src/Development/WeilBound/ThetaLifts/ConeHypotheses/CtSmallCone.lean}{115}{nonzero denominator in the primed tower}.

In the positive primed Witt index case descent gives a nonzero constant-term component. If its theta integral were zero, the zero pair would also split the constant term, contradicting uniqueness of splitting. For \(q=0\) use the supplied bottom nonvanishing. Choose the resulting right-hand polynomial test \(\Phi'_0\) and a point \(h_0\) with \(I_{V'}(\Phi'_0)(h_0)\ne0\). By the matching-existence hypothesis choose a left-hand match \(\Phi_0\) for this particular test. Continuity extends the polynomial comparison to every matched limit with the same \(c\). If \(c'\) is another comparison constant, subtract the two identities for \((\Phi_0,\Phi'_0,h_0)\). This gives \((c-c')I_{V'}(\Phi'_0)(h_0)=0\), hence \(c=c'\). \src{src/Development/WeilBound/ThetaLifts/ConeHypotheses/CtSmallCone.lean}{138}{Nonzero test integral}, \src{src/Development/WeilBound/ThetaLifts/ConeHypotheses/CtSmallCone.lean}{179}{passage to the matched limit} and \src{src/Development/WeilBound/ThetaLifts/ConeHypotheses/CtSmallCone.lean}{192}{matching the nonzero right-hand test}.
\end{proof}

The real and finite matching maps come from the spherical closure with real Lie-algebra lifting and the split, inert, and ramified local models. The strong identities used in the preceding proof are filled, in particular, by \code{eisBL} and the finite standard Siegel--Weil formula at the end of the spherical reduction. The proof above retains those identities and their nonvanishing hypotheses at their supplied scope, including the specific descent value of the \(L\)-factor. The actual assembly is in \src{src/Development/WeilBound/ThetaLifts/Reductions/FiniteClosure/TheoremKOfMod.lean}{14}{local models in the spherical closure}, \src{src/Development/WeilBound/ThetaLifts/SphericalLattices/SphericalClosure/CloseSph.lean}{188}{terminal comparison}, and \src{src/Development/AutomorphicArithmetic/ThetaCorrespondence/ThetaPairLifts/ThirdPlaceSeries/EveryCharacterOfLiteratureOfVWin.lean}{561}{supply of the two identities}. At the main consumer the integral comparison and local matching define the evaluation grid. Their transport to the square model is used there. See \src{src/Development/AutomorphicArithmetic/ThetaCorrespondence/ThetaPairLifts/ThirdPlaceSeries/EveryCharacterOfLiteratureOfVWin.lean}{810}{active theta consumer} and \src{src/Development/AutomorphicArithmetic/ThetaCorrespondence/ThetaPairLifts/ThirdPlaceSeries/SiegelEisensteinSecondTermHypotheses.lean}{50}{grid and associated bundles}.

\subsection{Component Hilbert--Samuel estimates and the final covolume}

For an admissible weight on a good deep cover write \(h_Z\) for the height of a component \(Z\), \(d_Z>0\) for its geometric degree, \(T\) for the section degree parameter and \(D_f\ge1\) for the flag dimension. The component Hilbert--Samuel input is \code{HSClaimS} with constant \(26\). It provides a lower degree bound through the auxiliary section parameters \(n,\Lambda,\beta,\Lambda_D,\Lambda_T,J\). The active numerical theorem arranges this as a lower bound on the degree in the associated bundle, with
\[
 B_Z=\binom{T-1}{D_f}d_Z,\quad
 z=7(D_f^2/T+D_f\rho),\quad
 A=26(26D_f^2+2zT),\quad
 R=4000b+2c.
\]
Here \(c\ge0\) is the cost per rank of the bundle transfer. Uniform choices are taken for the same point and stage across the components. See \src{src/Development/AutomorphicArithmetic/ThetaCorrespondence/ThetaPairLifts/ThirdPlaceSeries/HilbertSamuelLowerBoundOfComponents.lean}{89}{uniform choices and component input}.

\begin{lemma}[Rank conversion in the component degree estimate]\label{geo:component}
Let \(h\in\mathbb R\), \(b,B,d,c\) be nonnegative, \(0\le z\le1/2\), and \((1-z)d\le B\). Suppose the component estimate, after rearranging the Hilbert--Samuel ratio, gives
\[
 \deg\ge B\bigl(Th-3(D_f+1)|h|-A\log(2+|h|)-J_0\bigr),
 \qquad J_0\le4000b,
\]
and \(\deg_M\ge\deg-dc\). Then
\[
 B\bigl(Th-3(D_f+1)|h|-A\log(2+|h|)-(4000b+2c)\bigr)
 \le\deg_M .
\]
\end{lemma}
\begin{proof}
The rank comparison and \(z\le1/2\) give \(d\le2B\). Since \(c\ge0\), the cost \(dc\) is at most \(2Bc\). Since \(B\ge0\), replacing \(J_0\) by \(4000b\) only decreases the degree lower bound. Combining the two comparisons proves the result, without multiplying any inequality by a possibly negative height. The argument includes \(B=0\), in which case \(d=0\).
\end{proof}

The numerical discharge in \src{src/Development/AbelianVarietyTheory/Proofs/DivisorProofs/DegBoundOfClaimNumbers.lean}{9}{exact numerical hypotheses} requires \(1\le P_0\), \(1\le b\), \(TP_0\le\sqrt b\), \(\log b\le P_0\), and \(0\le\rho\) in addition to the rank comparison. Squaring the nonnegative sides yields \((TP_0)^2\le b\), while the formula for \(z\) gives \(D_f^2\le T/14\). The theorem combines these with the logarithmic bounds for the section parameters to obtain \(J_0\le4000b\). It applies the rank conversion of Lemma~\ref{geo:component} at \src{src/Development/AbelianVarietyTheory/Proofs/DivisorProofs/DegBoundOfClaimNumbers.lean}{121}{last remainder and transfer bounds}. The analytic estimate producing the component \code{HSClaimS} and its auxiliary section parameters is a premise at this stage.

An independent height comparison supplies \(c_s,c_c,X\ge0\) and real component slopes \(\ell_Z\) satisfying
\[
 0\le h_Z+c_s,\qquad
 (1-\rho)\max(0,\ell_Z+c_s)\le h_Z+c_s+c_c,\qquad
 -X\le \frac{\sum_Z\ell_Z}{\#\{Z\}}.
\]
The quantities \(TX\) and \(Tc_c\) and the prescribed multiple
\(c_s(6D_f+10+T(z+2\rho))\) are each at most \(C_1 E_P\).
The short-section Bergman kernel produces the first two comparisons, and
\code{MeanAll} supplies the average comparison. The remaining hypotheses
are positive component degrees comparable to a common degree, the rank
shares, and the inequality
\[
 2\frac{\sum_Z\deg_C(Z)}{\sum_Z\mathrm{rk}_C(Z)}
 \le \mu_{\rm cover}.
\]
The component bound above contributes the uniform budget
\(A\log\max(A,1)+R\le C_2E_P\).
With these inputs, the covering-slope result used here is
\(\mu_{\rm cover}\ge-2(3C_1+C_2+4)E_P\). The subsequent descent moves this
bound to the bundles used on the grid. The geometric comparison and its
numerical aggregation are used at the stated scope. See \src{src/Development/AutomorphicArithmetic/ThetaCorrespondence/ThetaPairLifts/ThirdPlaceSeries/HeightLowerBoundGeneral.lean}{21}{separate height comparison}, \src{src/Development/AutomorphicArithmetic/ThetaCorrespondence/ThetaPairLifts/ThirdPlaceSeries/ULowOf.lean}{16}{covering slope estimate}, and \src{src/Development/AutomorphicArithmetic/ThetaCorrespondence/ThetaPairLifts/ThirdPlaceSeries/TopCovolG.lean}{49}{deep cover and descent}.

\begin{proposition}[Uniform exponential absorption]\label{geo:covolume}
Fix \(a>0\), \(A_0,B_0\in\mathbb R\), and \(C_1,C_2\ge0\). Suppose the pointwise Gram, recurrence, cut, bad-set, lattice and size bounds yield
\[
 \mu_{\rm type}\le\mu_{\rm good}
 +12e^{-x^a}(x^{A_0}+x^{B_0}+C_1+C_2+1)E_P
\]
with \(E_P\ge0\). Suppose the core comparison and low-slope bound give \(\mu_{\rm good}\le \mathrm{mainT}+(C_1+C_2)E_P\), and \(\mu_{\rm type}\) identifies with the normalized log-covolume of the selected lattice. Then for all sufficiently large \(x\) the normalized log-covolume is at most \(\mathrm{mainT}+(C_1+C_2+1)E_P\), uniformly over the points and stages satisfying the hypotheses with these constants.
\end{proposition}
\begin{proof}
Set \(y=x^a\). For each exponent \(A_0/a,B_0/a\) choose a positive integer \(n\) larger than both exponents and zero. For \(y\ge1\), \(e^y\ge y^n/n!\). Each \(e^{-y}y^{A_0/a}\), \(e^{-y}y^{B_0/a}\), and \(e^{-y}\) tends to zero. The coefficient of \(E_P\) in the first inequality is therefore at most one after a threshold depending only on the fixed constants. Multiply this estimate by \(E_P\ge0\) and substitute the upper bound on \(\mu_{\rm good}\). The case \(E_P=0\) has the stated conclusion at every point satisfying the premises.
\end{proof}

In the actual construction \(a=e_z/7\) and the pointwise inequalities are those proved and combined in \src{src/Development/AbelianVarietyTheory/Proofs/DivisorProofs/CovolumeFromGramAndBounds.lean}{39}{covolume from the cut and slope data}. The family argument takes the maximum of this numerical threshold and the geometric thresholds, with \(C_1+C_2+1\) as the minimum error constant. The \code{hX} used at the main entry is obtained through the Chow-height and deep-cover inputs in \src{src/Development/AutomorphicArithmetic/ThetaCorrespondence/ThetaPairLifts/ThirdPlaceSeries/TopCovolSlim2Use.lean}{16}{covolume input on the actual family}. It is applied to the associated bundles on the evaluation grid to give \code{L3}, the inequality \(\mathrm{lcv}(X_{W,i},M_i)/D(P)\le N.l3_i\). See \src{src/Development/AutomorphicArithmetic/ThetaCorrespondence/ThetaPairLifts/ThirdPlaceSeries/LinkedBundlesL3OfOpen.lean}{31}{putting the covolume bound into L3}. Thus theta matching is used to realize the grid, while the Hilbert--Samuel estimates supply the quantitative slope information for the lattices selected on it.
